
\documentclass[a4paper,11pt]{article}
\usepackage[a4paper, margin=8em]{geometry}

% usa i pacchetti per la scrittura in italiano
\usepackage[french,italian]{babel}
\usepackage[T1]{fontenc}
\usepackage[utf8]{inputenc}
\frenchspacing 

% usa i pacchetti per la formattazione matematica
\usepackage{amsmath, amssymb, amsthm, amsfonts}

% usa altri pacchetti
\usepackage{gensymb}
\usepackage{hyperref}
\usepackage{standalone}

\usepackage{colortbl}

\usepackage{xstring}
\usepackage{karnaugh-map}

% imposta il titolo
\title{Embedded Operating Systems and Architectures notes}
\author{Luca Seggiani}
\date{2026}

% imposta lo stile
% usa helvetica
\usepackage[scaled]{helvet}
% usa palatino
\usepackage{palatino}
% usa un font monospazio guardabile
\usepackage{lmodern}

\renewcommand{\rmdefault}{ppl}
\renewcommand{\sfdefault}{phv}
\renewcommand{\ttdefault}{lmtt}

% circuiti
\usepackage{circuitikz}
\usetikzlibrary{babel}

% testo cerchiato
\newcommand*\circled[1]{\tikz[baseline=(char.base)]{
            \node[shape=circle,draw,inner sep=2pt] (char) {#1};}}

% disponi il titolo
\makeatletter
\renewcommand{\maketitle} {
	\begin{center} 
		\begin{minipage}[t]{.8\textwidth}
			\textsf{\huge\bfseries \@title} 
		\end{minipage}%
		\begin{minipage}[t]{.2\textwidth}
			\raggedleft \vspace{-1.65em}
			\textsf{\small \@author} \vfill
			\textsf{\small \@date}
		\end{minipage}
		\par
	\end{center}

	\thispagestyle{empty}
	\pagestyle{fancy}
}
\makeatother

% disponi teoremi
\usepackage{tcolorbox}
\newtcolorbox[auto counter, number within=section]{theorem}[2][]{%
	colback=blue!10, 
	colframe=blue!40!black, 
	sharp corners=northwest,
	fonttitle=\sffamily\bfseries, 
	title=Theorem~\thetcbcounter: #2, 
	#1
}

% disponi definizioni
\newtcolorbox[auto counter, number within=section]{definition}[2][]{%
	colback=red!10,
	colframe=red!40!black,
	sharp corners=northwest,
	fonttitle=\sffamily\bfseries,
	title=Definition~\thetcbcounter: #2,
	#1
}

% disponi codice
\usepackage{listings}
\usepackage[table]{xcolor}

\definecolor{codegreen}{rgb}{0,0.6,0}
\definecolor{codegray}{rgb}{0.5,0.5,0.5}
\definecolor{codepurple}{rgb}{0.58,0,0.82}
\definecolor{backcolour}{rgb}{0.95,0.95,0.92}

\lstdefinestyle{codestyle}{
		backgroundcolor=\color{black!5}, 
		commentstyle=\color{codegreen},
		keywordstyle=\bfseries\color{magenta},
		numberstyle=\sffamily\tiny\color{black!60},
		stringstyle=\color{green!50!black},
		basicstyle=\ttfamily\footnotesize,
		breakatwhitespace=false,         
		breaklines=true,                 
		captionpos=b,                    
		keepspaces=true,                 
		numbers=left,                    
		numbersep=5pt,                  
		showspaces=false,                
		showstringspaces=false,
		showtabs=false,                  
		tabsize=2
}

\lstdefinestyle{shellstyle}{
		backgroundcolor=\color{black!5}, 
		basicstyle=\ttfamily\footnotesize\color{black}, 
		commentstyle=\color{black}, 
		keywordstyle=\color{black},
		numberstyle=\color{black!5},
		stringstyle=\color{black}, 
		showspaces=false,
		showstringspaces=false, 
		showtabs=false, 
		tabsize=2, 
		numbers=none, 
		breaklines=true
}

\lstdefinelanguage{x86}{ 
  keywords={AAA, AAD, AAM, AAS, ADC, ADCB, ADCW, ADCL, ADD, ADDB, ADDW, ADDL, AND, ANDB, ANDW, ANDL,
        ARPL, BOUND, BSF, BSFL, BSFW, BSR, BSRL, BSRW, BSWAP, BT, BTC, BTCB, BTCW, BTCL, BTR, 
        BTRB, BTRW, BTRL, BTS, BTSB, BTSW, BTSL, CALL, CBW, CDQ, CLC, CLD, CLI, CLTS, CMC, CMP,
        CMPB, CMPW, CMPL, CMPS, CMPSB, CMPSD, CMPSW, CMPXCHG, CMPXCHGB, CMPXCHGW, CMPXCHGL,
        CMPXCHG8B, CPUID, CWDE, DAA, DAS, DEC, DECB, DECW, DECL, DIV, DIVB, DIVW, DIVL, ENTER,
        HLT, IDIV, IDIVB, IDIVW, IDIVL, IMUL, IMULB, IMULW, IMULL, IN, INB, INW, INL, INC, INCB,
        INCW, INCL, INS, INSB, INSD, INSW, INT, INT3, INTO, INVD, INVLPG, IRET, IRETD, JA, JAE,
        JB, JBE, JC, JCXZ, JE, JECXZ, JG, JGE, JL, JLE, JMP, JNA, JNAE, JNB, JNBE, JNC, JNE, JNG,
        JNGE, JNL, JNLE, JNO, JNP, JNS, JNZ, JO, JP, JPE, JPO, JS, JZ, LAHF, LAR, LCALL, LDS,
        LEA, LEAVE, LES, LFS, LGDT, LGS, LIDT, LMSW, LOCK, LODSB, LODSD, LODSW, LOOP, LOOPE,
        LOOPNE, LSL, LSS, LTR, MOV, MOVB, MOVW, MOVL, MOVSB, MOVSD, MOVSW, MOVSX, MOVSXB,
        MOVSXW, MOVSXL, MOVZX, MOVZXB, MOVZXW, MOVZXL, MUL, MULB, MULW, MULL, NEG, NEGB, NEGW,
        NEGL, NOP, NOT, NOTB, NOTW, NOTL, OR, ORB, ORW, ORL, OUT, OUTB, OUTW, OUTL, OUTSB, OUTSD,
        OUTSW, POP, POPL, POPW, POPB, POPA, POPAD, POPF, POPFD, PUSH, PUSHL, PUSHW, PUSHB, PUSHA, 
        RCLL, RCR, RCRB, RCRW, RCRL, RDMSR, RDPMC, RDTSC, REP, REPE, REPNE, RET, ROL, ROLB, ROLW,
        ROLL, ROR, RORB, RORW, RORL, SAHF, SAL, SALB, SALW, SALL, SAR, SARB, SARW, SARL, SBB,
        SBBB, SBBW, SBBL, SCASB, SCASD, SCASW, SETA, SETAE, SETB, SETBE, SETC, SETE, SETG, SETGE,
        SETL, SETLE, SETNA, SETNAE, SETNB, SETNBE, SETNC, SETNE, SETNG, SETNGE, SETNL, SETNLE,
        SETNO, SETNP, SETNS, SETNZ, SETO, SETP, SETPE, SETPO, SETS, SETZ, SGDT, SHL, SHLB, SHLW,
        SHLL, SHLD, SHR, SHRB, SHRW, SHRL, SHRD, SIDT, SLDT, SMSW, STC, STD, STI, STOSB, STOSD,
        STOSW, STR, SUB, SUBB, SUBW, SUBL, TEST, TESTB, TESTW, TESTL, VERR, VERW, WAIT, WBINVD,
        XADD, XADDB, XADDW, XADDL, XCHG, XCHGB, XCHGW, XCHGL, XLAT, XLATB, XOR, XORB, XORW, XORL},
  keywordstyle=\color{blue}\bfseries,
  ndkeywordstyle=\color{darkgray}\bfseries,
  identifierstyle=\color{black},
  sensitive=false,
  comment=[l]{\#},
  morecomment=[s]{/*}{*/},
  commentstyle=\color{purple}\ttfamily,
  stringstyle=\color{red}\ttfamily,
  morestring=[b]',
  morestring=[b]"
}

\lstdefinelanguage{riscv}{
  keywords={
    LUI, AUIPC,
    JAL, JALR,
    BEQ, BNE, BLT, BGE, BLTU, BGEU,
    LB, LH, LW, LBU, LHU, LD,
    SB, SH, SW, SD,
    ADDI, SLTI, SLTIU, XORI, ORI, ANDI,
    SLLI, SRLI, SRAI,
    ADD, SUB, SLL, SLT, SLTU, XOR, SRL, SRA, OR, AND,
    ADDIW, SLLIW, SRLIW, SRAIW,
    ADDW, SUBW, SLLW, SRLW, SRAW,
    MUL, MULH, MULHSU, MULHU, DIV, DIVU, REM, REMU,
    MULW, DIVW, DIVUW, REMW, REMUW,
    LR.W, SC.W, AMOSWAP.W, AMOADD.W, AMOXOR.W,
    AMOAND.W, AMOOR.W, AMOMIN.W, AMOMAX.W,
    AMOMINU.W, AMOMAXU.W,
    LR.D, SC.D, AMOSWAP.D, AMOADD.D, AMOXOR.D,
    AMOAND.D, AMOOR.D, AMOMIN.D, AMOMAX.D,
    AMOMINU.D, AMOMAXU.D,
    FLW, FSW,
    FMADD.S, FMSUB.S, FNMSUB.S, FNMADD.S,
    FADD.S, FSUB.S, FMUL.S, FDIV.S, FSQRT.S,
    FSGNJ.S, FSGNJN.S, FSGNJX.S,
    FMIN.S, FMAX.S,
    FCVT.W.S, FCVT.WU.S, FCVT.S.W, FCVT.S.WU,
    FMV.X.W, FMV.W.X,
    FEQ.S, FLT.S, FLE.S,
    FLD, FSD,
    FMADD.D, FMSUB.D, FNMSUB.D, FNMADD.D,
    FADD.D, FSUB.D, FMUL.D, FDIV.D, FSQRT.D,
    FSGNJ.D, FSGNJN.D, FSGNJX.D,
    FMIN.D, FMAX.D,
    FCVT.W.D, FCVT.WU.D, FCVT.D.W, FCVT.D.WU,
    FCVT.S.D, FCVT.D.S,
    FMV.X.D, FMV.D.X,
    FEQ.D, FLT.D, FLE.D,
    CSRRW, CSRRS, CSRRC,
    CSRRWI, CSRRSI, CSRRCI,
    FENCE.I,
    FENCE,
    ECALL, EBREAK,
    MRET, SRET, URET,
    WFI,
    SFENCE.VMA,
    NOP, LI, LA, MV,
    NOT, NEG, NEGW,
    SEQZ, SNEZ, SLTZ, SGTZ,
    BEQZ, BNEZ, BLEZ, BGEZ, BLTZ, BGTZ,
    BGT, BLE, BGTU, BLEU,
    J, JR, RET,
    CALL, TAIL,
    RDINSTRET, RDCYCLE, RDTIME,
    CSRR, CSRW, CSRS, CSRC,
    CSRWI, CSRSI, CSRCI
  },
  keywordstyle=\color{blue}\bfseries,
  ndkeywordstyle=\color{darkgray}\bfseries,
  identifierstyle=\color{black},
  sensitive=false,
  comment=[l]{\#},
  morecomment=[s]{/*}{*/},
  commentstyle=\color{purple}\ttfamily,
  stringstyle=\color{red}\ttfamily,
  morestring=[b]',
  morestring=[b]"
}

\lstdefinelanguage{arm}{
  keywords={
    ADC, ADD, AND, BIC, CMN, CMP, EOR, MOV, MVN,
    RSB, RSC, SBC, SUB, TST, TEQ,
    MLA, MLS, MUL, SMLAL, SMULL, UMLAL, UMULL,
    QADD, QSUB, QDADD, QDSUB,
    QADD16, QADD8, QASX, QSAX,
    QSUB16, QSUB8, QSAX, QSUBADDX,
    ASR, LSL, LSR, ROR, RRX,
    B, BL, BX, BLX,
    BXJ,
    BEQ, BNE, BCS, BHS, BCC, BLO,
    BMI, BPL, BVS, BVC,
    BHI, BLS, BGE, BLT, BGT, BLE,
    BAL,
    LDR, LDRB, LDRH, LDRSB, LDRSH,
    LDRD, LDM, LDMIA, LDMIB, LDMDA, LDMDB,
    LDRT, LDRBT, LDRHT, LDRSBT, LDRSHT,
    STR, STRB, STRH, STRD,
    STM, STMIA, STMIB, STMDA, STMDB,
    STRT, STRBT, STRHT,
    PUSH, POP,
    SWP, SWPB,
    REV, REV16, REVSH,
    BFC, BFI, BFXIL, SBFX, UBFX,
    CLZ,
    LDREX, LDREXB, LDREXH, LDREXD,
    STREX, STREXB, STREXH, STREXD,
    DMB, DSB, ISB,
    MCR, MCRR, MRC, MRRC,
    MRS, MSR,
    CPS, SETEND, SVC, BKPT,
    CLREX, WFE, WFI, SEV,
    YIELD, NOP,
    VLDR, VSTR,
    VMOV, VABS, VNEG,
    VADD, VSUB, VMUL, VDIV,
    VMLA, VMLS, VNMLA, VNMLS, VNMUL,
    VMAX, VMIN,
    VSQRT,
    VCMP, VCMPE,
    VCVT,
    VCLZ, VCLS,
    VAND, VORR, VEOR, VBIC, VMVN,
    VLD1, VLD2, VLD3, VLD4,
    VST1, VST2, VST3, VST4,
    ADR, ADCS, ADDS, SUBS, MOVS, ANDS, ORRS, EORS
  },
  keywordstyle=\color{blue}\bfseries,
  ndkeywordstyle=\color{darkgray}\bfseries,
  identifierstyle=\color{black},
  sensitive=false,
  comment=[l]{@},
  morecomment=[l]{;},
  morecomment=[s]{/*}{*/},
  commentstyle=\color{purple}\ttfamily,
  stringstyle=\color{red}\ttfamily,
  morestring=[b]',
  morestring=[b]"
}

\lstset{style=codestyle, language=C}

% disponi sezioni
\usepackage{titlesec}

\titleformat{\section}
	{\sffamily\Large\bfseries} 
	{\thesection}{1em}{} 
\titleformat{\subsection}
	{\sffamily\large\bfseries}   
	{\thesubsection}{1em}{} 
\titleformat{\subsubsection}
	{\sffamily\normalsize\bfseries} 
	{\thesubsubsection}{1em}{}

% tikz
\usepackage{tikz}

% float
\usepackage{float}

% grafici
\usepackage{pgfplots}
\pgfplotsset{width=10cm,compat=1.9}

% disponi alberi
\usepackage{forest}

\forestset{
	rectstyle/.style={
		for tree={rectangle,draw,font=\large\sffamily}
	},
	roundstyle/.style={
		for tree={circle,draw,font=\large}
	}
}

% disponi algoritmi
\usepackage{algorithm}
\usepackage{algorithmic}
\makeatletter
\renewcommand{\ALG@name}{Algorithm}
\makeatother

% disponi numeri di pagina
\usepackage{fancyhdr}
\fancyhf{} 
\fancyfoot[L]{\sffamily{\thepage}}

\makeatletter
\fancyhead[L]{\raisebox{1ex}[0pt][0pt]{\sffamily{\@title \ \@date}}} 
\fancyhead[R]{\raisebox{1ex}[0pt][0pt]{\sffamily{\@author}}}
\makeatother

\begin{document}
% sezione (data)
\section{Lesson 07-10-26}

% stili pagina
\thispagestyle{empty}
\pagestyle{fancy}

% testo
\subsection{ARM hardware interface}

The ARM hardware interface is split in 2 components:
\begin{itemize}
	\item The \textbf{CMDSDK} (\textit{Cortex-M System Design Kit}): it is an hardware specification of ready-made peripheral designs that chip designers place next to the CPU;
	\item The \textbf{CMSIS} (\textit{Cortex-M Software Interface Standard}): it is a standard for the C headers and functions that describe a chip.
\end{itemize}

Our board's peripherals are CMSDK blocks: Arm documents their registers once, the same on every chip that
uses them.
More specifically, it is an \textit{MPS2}: an Arm development board built around an FPGA (a chip whose logic is loaded at power-on). \textit{AN385}
(Application Note 385) is one FPGA image for it: a Cortex-M3 plus CMSDK peripherals.

\begin{table}[h!]
	\center \rowcolors{2}{white}{black!10}
	\begin{tabular} { c | c }
		\bfseries Address & \bfseries Device \\
		\hline 
		\texttt{0x4000\_0000} / \texttt{0x4000\_1000} / \texttt{0x4000\_1000} & Timer 0, Timer 1, dual timer \\

		\texttt{0x4000\_4000} / \texttt{0x4000\_5000} / \texttt{0x4000\_6000} & UART0, UART1, UART2 \\
		\texttt{0x4000\_8000} / \texttt{0x4001\_0000} & Watchdog, GPIO pins \\
		\texttt{0x4002\_8000} & FPGA I/O (MPS2 block)
	\end{tabular}
\end{table}

We use no vendor CMSIS files (no libraries at all): our register header is hand-written in CMSIS style, so its
names look like those in vendor code.

Basically, we have a CMDSDK peripheral set, for which we program a CMSIS interface by hand.

\subsubsection{Buses}

A \textbf{bus} is the set of wires (address, data, control) that carries every load and store from the CPU to a memory or a
device. Arm defines two kinds:

\begin{itemize}
\item \textbf{AHB} (\textit{Advanced High-performance Bus}): the main bus; pipelined, one transfer per cycle. Memories and devices
that need speed sit here;
\item \textbf{APB} (\textit{Advanced Peripheral Bus}): simple, slow, low power. For registers that are touched rarely: configure a
timer, send one byte
\end{itemize}

The bridge turns an AHB access into an APB access: a few extra cycles are needed for conversion.
For our code, nothing changes: a register is a load or store at an address on either bus; only the timing differs.

\subsubsection{External interrupts}

A device can raise an interrupt request (\textbf{IRQ}) line to tell the CPU "something happened". Device line $N$ is served by vector-table slot $16 + N$, right after the 16 core slots we saw were related to internal interrupts (exceptions, such as the Reset Handler). 

\begin{table}[H]
\centering
\begin{tabular}{c|l@{\hspace{2cm}}c|l}
\hline
\textbf{IRQ} & \textbf{Interrupt source} & \textbf{IRQ} & \textbf{Interrupt source} \\
\hline
0  & UART 0 receive & 16 & GPIO 2 combined \\
1  & UART 0 transmit & 17 & GPIO 3 combined \\
2  & UART 1 receive & 18 & UART 3 receive \\
3  & UART 1 transmit & 19 & UART 3 transmit \\
4  & UART 2 receive & 20 & UART 4 receive \\
5  & UART 2 transmit & 21 & UART 4 transmit \\
6  & GPIO 0 combined & 22 & SPI 2 \\
7  & GPIO 1 combined & 23 & SPI 3, SPI 4 \\
8  & Timer 0 & 24 & GPIO 0, pin 0 \\
9  & Timer 1 & 25 & GPIO 0, pin 1 \\
10 & Dual Timer & 26 & GPIO 0, pin 2 \\
11 & SPI 0, SPI 1 & 27 & GPIO 0, pin 3 \\
12 & UART overflow (UART 0, 1, 2) & 28 & GPIO 0, pin 4 \\
13 & Ethernet & 29 & GPIO 0, pin 5 \\
14 & Audio I2S & 30 & GPIO 0, pin 6 \\
15 & Touch Screen Controller & 31 & GPIO 0, pin 7 \\
\hline
\end{tabular}
\end{table}

\subsection{ARM ABI}

Let's look at the ARM \textbf{ABI} (\textit{Application Binary Interface}), as it will be useful to know how to implement interrupt handlers (we have to know how arguments are passed around).

C and assembly can call each other only if both follow the same rules: the \textbf{AAPCS} (\textit{Procedure Call Standard for the Arm Architecture}).

\begin{center}
	\begin{tabular}{c|p{5cm}|p{5cm}}
\hline
\textbf{Register} & \textbf{Role in a call} & \textbf{Kept across the call?} \\
\hline
r0--r3 & The first four arguments; r0 also carries the result & No: \textbf{caller-saved} (scratch) \\
\hline
r4--r11 & The function's own long-lived values & Yes: \textbf{callee-saved} \\
\hline
r12 & Scratch & No: \textbf{caller-saved} \\
\hline
lr & Return address, written by \texttt{bl} & No: a function that calls another saves it first \\
\hline
sp & Must be a multiple of 8 at every call & Yes \\
\hline
\end{tabular}
\end{center}

\begin{itemize}
\item \textbf{Caller-saved}: if the caller still needs r0--r3 or r12 after the call, it saves them itself;
\item \textbf{Callee-saved}: a function that wants to modify r4 must restore it before returning.
\end{itemize}

\end{document}
